\chapter{Nominalstil}
\label{cha:Nominalstil}

Der Nominalstil (auch „Nominalisierung") ersetzt Verben durch Substantive. Er ist charakteristisch für den formellen, akademischen Stil.

\section{Verben $\rightarrow$ Nomen}
\label{sec:VerbZuNomen}

\begin{center}
\begin{tabular}{|l|l|}
\hline
\textbf{Verb} & \textbf{Nomen} \\ \hline
entscheiden & die Entscheidung \\ \hline
planen & die Planung \\ \hline
analysieren & die Analyse \\ \hline
verbessern & die Verbesserung \\ \hline
entwickeln & die Entwicklung \\ \hline
erforschen & die Erforschung \\ \hline
bearbeiten & die Bearbeitung \\ \hline
beantworten & die Beantwortung \\ \hline
\end{tabular}
\end{center}

\section{Übungen zum Nominalstil}
\label{sec:NominalstilUebungen}

\begin{enumerate}
    \item Wir planen die Reise. $\rightarrow$ Nominalstil \quad \textit{(Lösung: Die Planung der Reise)}
    \item Sie analysieren die Daten. $\rightarrow$ Nominalstil \quad \textit{(Lösung: Die Analyse der Daten)}
    \item Er verbessert die Software. $\rightarrow$ Nominalstil \quad \textit{(Lösung: Die Verbesserung der Software)}
    \item Wir entwickeln neue Technologien. $\rightarrow$ Nominalstil \quad \textit{(Lösung: Die Entwicklung neuer Technologien)}
    \item Er forscht auf diesem Gebiet. $\rightarrow$ Nominalstil \quad \textit{(Lösung: Die Erforschung auf diesem Gebiet)}
\end{enumerate}

\section{Stilistische Hinweise}
\label{sec:NominalstilHinweise}

\begin{tcolorbox}[title=Stil]
\begin{itemize}
    \item \textbf{Nominalstil}: formell, akademisch, schriftlich
    \item \textbf{Verbalstil}: informell, mündlich
\end{itemize}
\end{tcolorbox}

Der Nominalstil sollte im akademischen Schreiben verwendet werden, in der gesprochenen Sprache klingt er jedoch oft steif.

\chapterend